% !TEX root = thesis.tex
\chapter{Introduction}
\label{sec:introduction}

The rapid development of \gls{ai} has led to its increasing use in areas such as autonomous systems, robotics, healthcare, and \gls{iot} devices. Many of these applications require data to be processed directly at the edge, where decisions must be made with low latency and limited communication with cloud servers. However, traditional \gls{ai} models often require considerable computational resources, memory bandwidth, and energy, which makes their deployment on resource-constrained embedded devices challenging. These limitations have increased the demand for specialized hardware architectures that can provide efficient \gls{ai} processing while supporting real-time operation and adaptation to changing environments.

\section{Motivation and Problem Statement}
\label{subsec:motivation-and-problem-statement}

\Glspl{snn} are brain-inspired neural networks in which information is represented and communicated through discrete spike events. In contrast to conventional \glspl{ann}, their temporal state and binary event communication allow specialized hardware to exploit sparse activity and avoid some unnecessary computation and data movement. This processing model is attractive for low-latency and resource-constrained edge systems, provided that the hardware architecture preserves the event-driven characteristics of the network \cite{Carpegna2025}.

\Glspl{fpga} provide a suitable platform for exploring such architectures because they offer configurable parallelism, deterministic execution, and direct control over data flow, numerical representation, and memory organization. Neural operations can therefore be mapped to application-specific datapaths while the implementation remains modifiable during development. Spiker+ follows this approach by generating configurable VHDL accelerators for \gls{snn} inference from a high-level network description \cite{Carpegna2025}.

Many \gls{fpga}-based \gls{snn} accelerators, including the baseline Spiker+ flow considered in this work, focus on inference using parameters obtained through offline training. Their synaptic weights are initialized when the hardware is generated and remain fixed during normal execution. This approach separates training from deployment, but it prevents the implemented network from adapting its synaptic state directly on the target device. Adding online learning is not only a matter of implementing a learning equation: the required neuron state must be available at the correct time, the weight memory must support runtime writes, and every update must remain synchronized with the synaptic address being processed.

\Gls{sdsp} is well suited to this problem because a synaptic update is triggered by a presynaptic event and depends on information available locally at the synapse and postsynaptic neuron, including the membrane potential and a calcium-related estimate of recent firing activity \cite{Brader2007}. It avoids a global error-propagation datapath, but its integration still requires coordinated access to neuron state, bounded digital weights, and writable synaptic storage. The central problem addressed by this thesis is therefore how to introduce persistent \gls{sdsp}-based weight adaptation into the generated Spiker+ hierarchy without replacing its established inference schedule and neuron organization.

\section{Objective and Contributions}
\label{subsec:objective-and-contributions}

The objective of this thesis is to design, implement, and evaluate an \gls{fpga}-based \gls{snn} accelerator that supports \gls{sdsp} weight updates during network operation. The open-source Spiker+ framework is used as the inference baseline \cite{Carpegna2025}. The architecture generated by Spiker+ is extended while preserving its multi-cycle control, neuron datapaths, and recurrent output competition. When learning is enabled, the modified accelerator updates selected synaptic weights and retains them for subsequent input events. When learning is disabled, the same hardware remains available for inference without issuing new weight writes.

The implemented system also includes the hardware and software interfaces required for board-level operation. The complete accelerator is packaged as a Vivado \gls{ip} core and connected to the Zynq processing system through AXI4-Lite control, an AXI DMA input path, and an interrupt interface. A PYNQ runtime prepares the spike representation, configures inference or supervised learning, transfers complete samples from DDR, and retrieves the output spike counts and classification result. This system boundary allows the learning architecture to be evaluated on the PYNQ-Z1 rather than only through RTL simulation.

The main contributions of this thesis are:

\begin{itemize}
    \item a synthesizable digital \gls{sdsp} datapath that combines membrane-potential and calcium conditions with bounded, finite-width weight updates;
    \item the integration of ten local learning lanes, writable excitatory memory, and an address-aligned read--modify--write path into the generated Spiker+ layer while retaining an inference-only operating mode;
    \item the packaging and deployment of the complete accelerator in a PYNQ-Z1 system with processor control, DDR-to-stream data movement, runtime learning configuration, and observable transaction results; and
    \item multilevel verification and hardware evaluation demonstrating persistent teacher-guided weight adaptation, timing closure at 100~MHz, continuous execution of 180,000 training transactions, and an increase in strict MNIST validation accuracy from 5.9\% to 39.6\% over three epochs.
\end{itemize}

The reported accuracy is an intermediate result rather than a converged performance claim. The training experiment uses a single-layer network, one parameter configuration, and three epochs; the results were still changing at the final checkpoint. The evaluation therefore supports the functionality and integration of the online-learning architecture more strongly than it supports the classification performance of the selected configuration.

\Gls{dpr} was considered during the design process but was not implemented in the final system. The \gls{sdsp} module depends on algorithm-specific neuron state, memory organization, and update timing, and neither a second interface-compatible learning module nor a runtime-switching use case was available. A dedicated reconfigurable partition would consequently add implementation and verification complexity without providing useful functionality for the evaluated system. Chapter \ref{sec:design_methodology} documents this applicability assessment and the resulting scope decision.

\section{Thesis Organization}
\label{subsec:thesis-organization}

The remainder of this thesis is organized as follows. Chapter \ref{sec:background} introduces \glspl{snn}, the \gls{lif} neuron model, \gls{sdsp}, \gls{fpga}-based acceleration, the concepts required for the \gls{dpr} assessment, and the Spiker+ framework. Chapter \ref{sec:related_work} reviews digital online-learning processors and positions the proposed Spiker+ extension relative to existing \gls{sdsp}- and \gls{stdp}-based systems. Chapter \ref{sec:design_methodology} derives the online-learning architecture, describes the digital weight-update path and its integration into Spiker+, and assesses the applicability of \gls{dpr}. Chapter \ref{sec:implementation} presents the final RTL hierarchy, the packaged accelerator IP, the Vivado block design, the PYNQ runtime, and the verification infrastructure. Chapter \ref{sec:results_and_evaluation} reports the functional, implementation, runtime, and three-epoch learning results and compares them with TEDOP. Finally, Chapter \ref{sec:conclusion_and_future_work} summarizes the findings and identifies directions for improving training, encoding, network structure, and measurement.
